\manualchapter{Tuning and modulation}{sec:modulation}
The frequency knob spans five octaves either side of middle C; its default
is C4. V/OCT adds one octave per volt. The multiplier spans 0--15, defaults
to 1, and uses $\tfrac12$ for 0. Start with integer ratios for pitched FM;
small frequency offsets produce beating or inharmonic sidebands.

\begin{table}[ht]
\centering\small
\begin{tabular}{@{}lrrrrrrrr@{}}
\toprule LFO index & 0 & 1 & 2 & 3 & 4 & 5 & 6 & 7 \\
Rate label (Hz) & 3.98 & 5.56 & 6.02 & 6.37 & 6.88 & 9.63 & 48.1 & 72.2 \\
\bottomrule
\end{tabular}
\caption{The eight internal LFO rate labels; default index 0.}
\label{tab:lfo-frequencies}
\end{table}

AMS selects amplitude modulation sensitivity: indices 0--3 correspond to
0, 1.4, 5.9 and 11.8\,dB. FMS selects frequency modulation sensitivity:
indices 0--7 correspond to 0, 3.4, 6.7, 10, 14, 20, 40 and 80 percent of a
semitone. Both default to 0, so changing the LFO rate alone is inaudible.

\subsection{Feedback, external FM and level}
Feedback spans 0--7 and defaults to 0. It feeds the operator back into
itself. There is no separate feedback CV input. The bipolar FM depth knob
spans $-1$ to $+1$ and defaults to zero. It scales the FM input, normalized
by 5\,V and limited to $\pm1$, before conversion to the operator's signed
modulation signal. This input modulates the operator, rather than adding
one octave per volt to the tuning control.

Output volume spans 0--127 and defaults to 127. VOL adds $127V/10$ to the
knob and clamps the result to 0--127. To open it with a positive envelope,
start the knob at zero; at its default a positive CV is already at the ceiling.
OUT is the chip's clipped 14-bit signal scaled to approximately $\pm5$\,V.
